\documentclass[12pt, a4paper]{article}

% ---------- packages ----------
\usepackage[T1]{fontenc}
\usepackage[utf8]{inputenc}

\usepackage[margin=2.5cm]{geometry}
\usepackage{amsmath, amsthm, amssymb}
\usepackage{mathtools}          % \coloneqq, \prescript, dcases
\usepackage{graphicx}
\usepackage{xcolor}
\usepackage{booktabs}
\usepackage{array}
\usepackage{hyperref}
\usepackage{tcolorbox}
\usepackage{mdframed}
\usepackage{parskip}            % space between paragraphs instead of indent
\usepackage{enumitem}
\usepackage{titlesec}

% ---------- colour palette ----------
\definecolor{accentblue}{RGB}{30, 60, 110}
\definecolor{lightblue}{RGB}{220, 232, 248}
\definecolor{warmgray}{RGB}{245, 242, 237}
\definecolor{borderblue}{RGB}{80, 120, 180}
\definecolor{notegreen}{RGB}{0, 100, 60}
\definecolor{darkorange}{RGB}{170, 80, 0}

% ---------- section formatting ----------
\titleformat{\section}{\large\bfseries\color{accentblue}}{}{0em}{\thesection\quad}[\vspace{-0.3em}\textcolor{accentblue}{\rule{\linewidth}{0.8pt}}]
\titleformat{\subsection}{\normalsize\bfseries\color{accentblue}}{}{0em}{\thesubsection\quad}

% ---------- tcolorbox styles ----------
\tcbuselibrary{breakable, skins}

\newtcolorbox{keybox}[1][]{
  enhanced, breakable,
  colback=lightblue, colframe=accentblue,
  boxrule=1pt, arc=4pt,
  fonttitle=\bfseries\small\color{accentblue},
  title={#1}
}

\newtcolorbox{resultbox}[1][]{
  enhanced, breakable,
  colback=warmgray, colframe=accentblue,
  boxrule=1.5pt, arc=4pt,
  fonttitle=\bfseries\small\color{accentblue},
  title={#1}
}

\newtcolorbox{notebox}[1][]{
  enhanced, breakable,
  colback=white, colframe=darkorange,
  boxrule=0.8pt, leftrule=4pt, arc=2pt,
  fonttitle=\bfseries\small\color{darkorange},
  title={#1}
}

% ---------- custom commands ----------
\newcommand{\E}{\mathbb{E}}
\newcommand{\Var}{\operatorname{Var}}
\newcommand{\dd}{\,\mathrm{d}}
\newcommand{\kap}{\kappa}

% label for numbered equations that are referenced
\newcommand{\ibpref}{\eqref{eq:ibp}}

% ---------- title ----------
\title{%
  \textbf{\textcolor{accentblue}{Continuous-Time Almgren--Chriss Model}}\\[4pt]
  \large Implementation Shortfall, Target Close, and TWAP\\
  \normalsize Static Mean--Variance Strategies, No Permanent Impact\\[6pt]
}
\author{}
\date{}

% ============================================================
\begin{document}
\maketitle
\thispagestyle{empty}
\tableofcontents
\newpage

% -------------------------------------------------------------------
\section{Target Close and TWAP orders}
\subsection{Present the continuous Almgren--Chriss model and find the optimal static strategy of IS, TC and TWAP orders in the mean--variance framework}
\subsubsection{Model Setup}
% -------------------------------------------------------------------

We set up the continuous-time model so we can derive the optimal liquidation schedules.

We liquidate a block of $Q > 0$ shares over a fixed horizon $[0,T]$.
We use the following quantities:

\begin{itemize}[leftmargin=1.5em]
  \item $q_t$: the number of shares still held (\emph{inventory}) at time $t$.
  \item $v_t = -\dot q_t$: the \emph{selling speed}, i.e.\ how many shares we sell per unit of time.
    A positive $v_t$ means we are selling.
  \item Boundary conditions: $q_0 = Q$ (we start with everything) and $q_T = 0$ (we finish flat). These conditions imply
        \[
          \int_0^T v_t \dd t = Q.
        \]
\end{itemize}

\paragraph{Price dynamics (no permanent impact).}
The unaffected mid-price follows a Brownian motion with no drift:
\[
  \dd S_t = \sigma \dd W_t, \qquad S_t = S_0 + \sigma W_t,
\]
where $\sigma > 0$ is the volatility and $W_t$ is a standard Brownian motion.

\paragraph{Temporary market impact.}
When we sell at speed $v_t$, we push the price against ourselves momentarily.
The actual execution price is
\[
  \tilde S_t = S_t - \eta v_t,
\]
where $\eta > 0$ is the temporary impact parameter.
Selling $v_t \dd t$ shares at this price generates cash:
\[
  \dd X_t = \tilde S_t \, v_t \dd t = (S_t - \eta v_t)\,v_t \dd t.
\]

\paragraph{Total cash received.}
Integrating over the whole horizon:
\[
  X_T = \int_0^T (S_t - \eta v_t)\,v_t \dd t.
\]

\paragraph{Criterion.}
We use a \emph{mean--variance} objective. For risk-aversion parameter $\lambda > 0$, we minimise
\[
  J = -\E[\Pi] + \lambda\,\Var(\Pi),
\]
where $\Pi$ is the \emph{payoff} of the strategy relative to a chosen benchmark.
We assume the selling schedule is fixed in advance and does not react to price moves.
Therefore, $q_t$ and $v_t$ are deterministic functions of time.

% -------------------------------------------------------------------
\subsubsection{A Shared Tool: Integration by Parts on the Cash Flow}
\label{sec:ibp}
% -------------------------------------------------------------------

We rewrite the cash flow in a form that we use in all three benchmark problems.

\paragraph{Expanding the cash flow.}
Plugging in $S_t = S_0 + \sigma W_t$ gives
\begin{align}
  X_T
  &= \int_0^T S_t\,v_t \dd t - \eta\int_0^T v_t^2 \dd t \notag\\
  &= S_0\!\int_0^T v_t \dd t
     + \sigma\!\int_0^T W_t\,v_t \dd t
     - \eta\!\int_0^T v_t^2 \dd t.
\end{align}
Using $\int_0^T v_t \dd t = Q$, the first term becomes $S_0 Q$:
\begin{equation}
  X_T = S_0 Q + \sigma\!\int_0^T W_t\,v_t \dd t - \eta\!\int_0^T v_t^2 \dd t.
  \label{eq:XT-raw}
\end{equation}

\paragraph{Stochastic integration by parts.}
The term $\int_0^T W_t \,v_t \dd t$ is an ordinary integral because $v_t$ is deterministic.
It is still random because it depends on $W_t$.
We rewrite it as a stochastic integral using integration by parts.
Since $v_t = -\dot q_t$, we have $v_t \dd t = -\dd q_t$.
We use
\[
  u = W_t, \quad \dd V = v_t \dd t = -\dd q_t:
\]
\[
  \int_0^T W_t\,v_t \dd t
  = -\int_0^T W_t \dd q_t
  = -\bigl[W_t\,q_t\bigr]_0^T + \int_0^T q_t \dd W_t.
\]
Evaluating the boundary term: $W_0 = 0$ and $q_T = 0$, so
\[
  \bigl[W_t\,q_t\bigr]_0^T = W_T\cdot 0 - 0\cdot Q = 0.
\]
\begin{equation}
  \boxed{\int_0^T W_t\,v_t \dd t = \int_0^T q_t \dd W_t.}
  \label{eq:ibp}
\end{equation}

\paragraph{Canonical form of the cash flow.}
Plugging \eqref{eq:ibp} into \eqref{eq:XT-raw} gives
\begin{equation}
  X_T = S_0 Q - \eta\!\int_0^T v_t^2 \dd t + \sigma\!\int_0^T q_t \dd W_t.
  \label{eq:XT-canonical}
\end{equation}
This form holds for any strategy.
We use it in each benchmark case below.

% -------------------------------------------------------------------
\subsubsection{Implementation Shortfall (IS)}
\label{sec:IS}
% -------------------------------------------------------------------

We solve the IS problem so we can identify the optimal schedule relative to the initial price.

\paragraph{Benchmark and payoff}

The \emph{implementation shortfall} benchmark is the initial price $S_0$.
The ideal outcome is to sell all $Q$ shares instantly at $S_0$ and receive $QS_0$.
The shortfall is the gap between that ideal value and the cash we actually receive:
\[
  \Pi_{\mathrm{IS}} = X_T - QS_0.
\]
Minimising the shortfall in the mean--variance sense is equivalent to maximising the cash $X_T$.
So we use the objective
\[
  J_{\mathrm{IS}} = -\E[\Pi_{\mathrm{IS}}] + \lambda\,\Var(\Pi_{\mathrm{IS}}).
\]
Since $QS_0$ is constant, it drops out of both the mean and the variance:
$\E[\Pi_{\mathrm{IS}}] = \E[X_T] - QS_0$ and $\Var(\Pi_{\mathrm{IS}}) = \Var(X_T)$.

\paragraph{Mean of \texorpdfstring{$X_T$}{XT}}

In \eqref{eq:XT-canonical}, the only random term is $\sigma\int_0^T q_t \dd W_t$.
Since $q_t$ is deterministic, the stochastic integral has zero expectation.
Therefore,
\[
  \E[X_T] = S_0 Q - \eta\!\int_0^T v_t^2 \dd t.
\]
We can drop the constant term $-S_0 Q$ from the optimisation.
This leaves
\[
  -\E[\Pi_{\mathrm{IS}}] = \eta\!\int_0^T v_t^2 \dd t = \eta\!\int_0^T \dot q_t^2 \dd t.
\]

\paragraph{Variance of \texorpdfstring{$X_T$}{XT}}

Only the stochastic integral contributes to the variance.
By It\^{o}'s isometry,
\[
  \Var(X_T)
  = \sigma^2\,\Var\!\left(\int_0^T q_t \dd W_t\right)
  = \sigma^2\!\int_0^T q_t^2 \dd t.
\]

\paragraph{Optimisation problem}

Putting the mean and variance together gives
\begin{equation}
  J_{\mathrm{IS}}[q]
  = \eta\!\int_0^T \dot q_t^2 \dd t + \lambda\sigma^2\!\int_0^T q_t^2 \dd t,
  \label{eq:JIS}
\end{equation}
subject to $q_0 = Q$, $q_T = 0$.

\begin{notebox}[Remark: Reading the objective]
  The first term $\eta\int\dot q_t^2$ is the \emph{impact cost}: selling faster
  (larger $|\dot q_t|$) costs more. The second term $\lambda\sigma^2\int q_t^2$
  is the \emph{market risk}: holding more inventory (larger $|q_t|$) exposes us to
  more price uncertainty. The optimal schedule balances these two forces.
\end{notebox}

\paragraph{Euler--Lagrange equation}

The integrand in \eqref{eq:JIS} is $\mathcal{L}(q, \dot q) = \eta\dot q^2 + \lambda\sigma^2 q^2$.
Applying the Euler--Lagrange equation
$\frac{d}{dt}\!\left(\frac{\partial\mathcal{L}}{\partial\dot q}\right)
 - \frac{\partial\mathcal{L}}{\partial q} = 0$
gives
\[
  \frac{d}{dt}(2\eta\dot q) - 2\lambda\sigma^2 q = 0
  \implies
  \ddot q - \kap^2 q = 0,
  \qquad \kap^2 \coloneqq \frac{\lambda\sigma^2}{\eta}.
\]

\paragraph{Solution}

The general solution of $\ddot q - \kap^2 q = 0$ is
\[
  q(t) = A\cosh(\kap t) + B\sinh(\kap t).
\]
Using $q(0) = Q$ gives $A = Q$.
Using $q(T) = 0$ gives
$B = -Q\dfrac{\cosh(\kap T)}{\sinh(\kap T)} = -Q\coth(\kap T)$.
$q(t)$ is therefore $Q\cosh(\kap t) - Q\coth(\kap T)\sinh(\kap t)$.
Using the hyperbolic identity
$\sinh(\kap(T-t)) = \sinh(\kap T)\cosh(\kap t) - \cosh(\kap T)\sinh(\kap t)$,
this gives the optimal IS schedule:

\begin{resultbox}[IS - optimal schedule]
\[
  q_{\mathrm{IS}}(t) = Q\,\frac{\sinh\!\bigl(\kap(T-t)\bigr)}{\sinh(\kap T)},
  \qquad
  v_{\mathrm{IS}}(t) = Q\,\frac{\kap\cosh\!\bigl(\kap(T-t)\bigr)}{\sinh(\kap T)},
  \qquad \kap^2 = \frac{\lambda\sigma^2}{\eta}.
\]
For large $\kap$ (high risk aversion or high volatility), the schedule front-loads
selling: $q_t$ drops quickly then flattens. For small $\kap$, it is close to linear.
\end{resultbox}

% -------------------------------------------------------------------
\subsubsection{Target Close (TC)}
\label{sec:TC}
% -------------------------------------------------------------------

We solve the TC problem so we can identify the optimal schedule relative to the close.

\paragraph{Benchmark and payoff}

The \emph{target close} benchmark is the terminal price $S_T$.
We measure performance relative to the price at the end of the trading window.
The payoff is
\[
  \Pi_{\mathrm{TC}} = X_T - QS_T.
\]

\paragraph{Simplifying the payoff}

We start from \eqref{eq:XT-canonical} and write $S_T = S_0 + \sigma W_T$:
\begin{align*}
  \Pi_{\mathrm{TC}}
  &= S_0 Q - \eta\!\int_0^T v_t^2 \dd t + \sigma\!\int_0^T q_t \dd W_t
     - Q(S_0 + \sigma W_T).
\end{align*}
The $S_0 Q$ terms cancel.
We write $W_T = \int_0^T \dd W_t$, so $QW_T = Q\int_0^T \dd W_t$:
\begin{align}
  \Pi_{\mathrm{TC}}
  &= - \eta\!\int_0^T v_t^2 \dd t + \sigma\!\int_0^T q_t \dd W_t
     - Q\sigma\!\int_0^T \dd W_t \notag\\
  &= - \eta\!\int_0^T v_t^2 \dd t + \sigma\!\int_0^T (q_t - Q)\,\dd W_t.
  \label{eq:PiTC}
\end{align}

\paragraph{Mean and variance}

\textbf{Mean.} The stochastic integral has zero expectation:
\[
  -\E[\Pi_{\mathrm{TC}}] = \eta\!\int_0^T v_t^2 \dd t = \eta\!\int_0^T \dot q_t^2 \dd t.
\]

\textbf{Variance.} By It\^{o}'s isometry:
\[
  \Var(\Pi_{\mathrm{TC}})
  = \sigma^2\!\int_0^T (q_t - Q)^2 \dd t.
\]

\paragraph{Optimisation problem}

\begin{equation}
  J_{\mathrm{TC}}[q]
  = \eta\!\int_0^T \dot q_t^2 \dd t + \lambda\sigma^2\!\int_0^T (q_t - Q)^2 \dd t,
  \label{eq:JTC}
\end{equation}
subject to $q_0 = Q$, $q_T = 0$.

\paragraph{Euler--Lagrange equation}

The integrand in \eqref{eq:JTC} is $\mathcal{L}_{\mathrm{TC}}(q, \dot q) = \eta\dot q^2 + \lambda\sigma^2(q - Q)^2$.
The partial derivatives are
\[
  \frac{\partial\mathcal{L}_{\mathrm{TC}}}{\partial\dot q} = 2\eta\dot q,
  \qquad
  \frac{\partial\mathcal{L}_{\mathrm{TC}}}{\partial q} = 2\lambda\sigma^2(q - Q).
\]
The Euler--Lagrange equation gives
\[
  \ddot q - \kap^2(q - Q) = 0, \qquad \kap^2 = \frac{\lambda\sigma^2}{\eta}.
\]

\paragraph{Change of variable}

We write $y(t) = q(t) - Q$, where $y(t)$ is the gap from a full position.
Then $\ddot y = \ddot q$, and the ODE becomes
\[
  \ddot y - \kap^2 y = 0.
\]
The boundary conditions become
\[
  y(0) = q(0) - Q = Q - Q = 0, \qquad y(T) = q(T) - Q = 0 - Q = -Q.
\]

\paragraph{Solution}

The general solution is $y(t) = A\cosh(\kap t) + B\sinh(\kap t)$.
Using $y(0) = 0$ gives $A = 0$.
Using $y(T) = -Q$ gives $B = -Q/\sinh(\kap T)$.
With $y(t) = -Q\sinh(\kap t)/\sinh(\kap T)$, this gives the optimal TC schedule:

\begin{resultbox}[TC - optimal schedule]
\[
  q_{\mathrm{TC}}(t) = Q\!\left(1 - \frac{\sinh(\kap t)}{\sinh(\kap T)}\right),
  \qquad
  v_{\mathrm{TC}}(t) = Q\,\frac{\kap\cosh(\kap t)}{\sinh(\kap T)},
  \qquad \kap^2 = \frac{\lambda\sigma^2}{\eta}.
\]
The schedule starts slowly (holding most of the position)
and accelerates toward the end --- the mirror image of IS.
\end{resultbox}

% -------------------------------------------------------------------
\subsubsection{TWAP Benchmark}
\label{sec:TWAP}
% -------------------------------------------------------------------

We solve the TWAP problem so we can identify the schedule that stays aligned with the benchmark path.

\paragraph{Benchmark and payoff}

The \emph{time-weighted average price} benchmark is the time-average of the mid-price:
\[
  B_{\mathrm{TWAP}} = \frac{1}{T}\int_0^T S_t \dd t.
\]
The payoff is
\[
  \Pi_{\mathrm{TWAP}} = X_T - Q \cdot B_{\mathrm{TWAP}}
  = X_T - \frac{Q}{T}\int_0^T S_t \dd t.
\]

\paragraph{Handling the benchmark term}

We simplify $\frac{Q}{T}\int_0^T S_t \dd t$ by plugging in $S_t = S_0 + \sigma W_t$:
\[
  \frac{Q}{T}\int_0^T S_t \dd t
  = \frac{Q}{T}\!\int_0^T S_0 \dd t + \frac{Q\sigma}{T}\!\int_0^T W_t \dd t
  = QS_0 + \frac{Q\sigma}{T}\!\int_0^T W_t \dd t.
\]

We rewrite the remaining integral $\int_0^T W_t \dd t$ with integration by parts.
We use $u = W_t$ and $\dd V = \dd t$, so $\dd u = \dd W_t$ and $V = t$:
\begin{align}
  \int_0^T W_t \dd t
  &= \Bigl[t\, W_t\Bigr]_0^T - \int_0^T t\, \dd W_t
   = T W_T - \int_0^T t\, \dd W_t.
   \label{eq:ibp-twap}
\end{align}
Writing $W_T = \int_0^T \dd W_t$, the first term becomes
$T W_T = T\int_0^T \dd W_t = \int_0^T T\,\dd W_t$.
So the two terms combine into a single stochastic integral:
\[
  \int_0^T W_t \dd t
  = \int_0^T (T - t)\,\dd W_t.
\]

So the benchmark term becomes
\[
  \frac{Q}{T}\int_0^T S_t \dd t
  = QS_0 + \frac{Q\sigma}{T}\int_0^T (T-t)\,\dd W_t.
\]

\begin{keybox}[Key point: Natural benchmark path]
  Define the \emph{deterministic linear schedule}
  \[
    f(t) \coloneqq Q\,\frac{T-t}{T},
  \]
  which starts at $Q$ when $t=0$ and reaches $0$ at $t=T$:
  it is the inventory path of a trader selling at a perfectly uniform rate $Q/T$.
  Equation \eqref{eq:PiTWAP} shows that the stochastic risk of a TWAP order
  is entirely governed by the \emph{deviation} $q_t - f(t)$.
  This requires no time-average of $q$ — the benchmark enters directly as a path.
\end{keybox}


\paragraph{Mean and variance}

\textbf{Mean.} The stochastic integral has zero expectation:
\[
  -\E[\Pi_{\mathrm{TWAP}}] = \eta\!\int_0^T v_t^2 \dd t = \eta\!\int_0^T \dot q_t^2 \dd t.
\]

\textbf{Variance.} By It\^{o}'s isometry applied to \eqref{eq:PiTWAP}:
\[
  \Var(\Pi_{\mathrm{TWAP}})
  = \sigma^2\!\int_0^T \!\left(q_t - Q\,\frac{T-t}{T}\right)^{\!2} \dd t.
\]

\paragraph{Optimisation problem}

\begin{equation}
  J_{\mathrm{TWAP}}[q]
  = \eta\!\int_0^T \dot q_t^2 \dd t
    + \lambda\sigma^2\!\int_0^T \!\left(q_t - Q\,\frac{T-t}{T}\right)^{\!2} \dd t,
  \label{eq:JTWAP}
\end{equation}
subject to $q_0 = Q$, $q_T = 0$.

\begin{notebox}[Remark: IS, TC, and TWAP]
The three objectives have identical structure:
\[
  J[q] = \eta\!\int_0^T\dot q_t^2\,\dd t + \lambda\sigma^2\!\int_0^T\bigl(q_t - f(t)\bigr)^2\dd t,
\]
where the ``target path'' $f$ changes by benchmark:
\begin{center}
  \renewcommand{\arraystretch}{1.4}
  \begin{tabular}{@{}lll@{}}
    \toprule
    Benchmark & Target path $f(t)$ & Intuition \\
    \midrule
    IS   & $0$              & Hold nothing; sell everything up front. \\
    TC   & $Q$              & Hold everything; sell everything at the end. \\
    TWAP & $Q(T-t)/T$  & Hold and sell at a constant linear rate. \\
    \bottomrule
  \end{tabular}
\end{center}
\end{notebox}

\paragraph{Euler--Lagrange equation}

The integrand in \eqref{eq:JTWAP} is $\mathcal{L}_{\mathrm{TWAP}}(q,\dot q, t) = \eta\dot q^2 + \lambda\sigma^2(q - f(t))^2$,
where $f(t) = Q(T-t)/T$ is a known deterministic function of time.
The partial derivatives are
\[
  \frac{\partial\mathcal{L}_{\mathrm{TWAP}}}{\partial\dot q} = 2\eta\dot q,
  \qquad
  \frac{\partial\mathcal{L}_{\mathrm{TWAP}}}{\partial q} = 2\lambda\sigma^2(q - f(t)).
\]
The Euler--Lagrange equation gives
\[
  \ddot q - \kap^2\!\left(q - Q\,\frac{T-t}{T}\right) = 0, \qquad \kap^2 = \frac{\lambda\sigma^2}{\eta}.
\]

\paragraph{Change of variable}

We write
\[
  y(t) \coloneqq q(t) - f(t) = q(t) - Q\,\frac{T-t}{T}.
\]
where $y(t)$ is the deviation from the linear schedule.
Since $f$ is linear, its second derivative vanishes: $\ddot f(t) = 0$.
So $\ddot y = \ddot q$, and the ODE becomes
\[
  \ddot y - \kap^2 y = 0.
\]
This is the same ODE as before, but the boundary conditions are now:
\[
  y(0) = q(0) - f(0) = Q - Q = 0, \qquad y(T) = q(T) - f(T) = 0 - 0 = 0.
\]
Both boundary values are zero.

\paragraph{Solution}

The general solution is $y(t) = A\cosh(\kap t) + B\sinh(\kap t)$.
Using $y(0) = 0$ gives $A = 0$.
Using $y(T) = 0$ gives $B\sinh(\kap T) = 0$.
Since $\kap > 0$, we have $\sinh(\kap T) > 0$, so $B = 0$.
Since $y(t) = 0$ for all $t \in [0, T]$, this gives the optimal TWAP schedule:

\begin{resultbox}[TWAP - optimal schedule]
\[
  q_{\mathrm{TWAP}}(t) = Q\!\left(1 - \frac{t}{T}\right),
  \qquad
  v_{\mathrm{TWAP}}(t) = \frac{Q}{T}.
\]
The optimal strategy is to sell at a \textbf{perfectly constant rate} $Q/T$
for all $t$, for \textbf{any} value of $\lambda > 0$ and any $\kap$.
This is an exact result, not an approximation.
\end{resultbox}

% -------------------------------------------------------------------
\subsubsection{Summary and plots}
\label{sec:summary}
% -------------------------------------------------------------------

We collect the three schedules here so we can compare them directly:

\begin{resultbox}[Summary - optimal schedules]
All three strategies share the same Euler--Lagrange structure,
$\ddot y - \kap^2 y = 0$ on the deviation $y(t) = q(t) - f(t)$,
with $\kap^2 = \lambda\sigma^2/\eta$.

\medskip
\begin{center}
  \renewcommand{\arraystretch}{2.0}
  \begin{tabular}{@{}llll@{}}
    \toprule
    \textbf{Benchmark} & \textbf{Optimal inventory} $q(t)$ & \textbf{BCs on $y$} & \textbf{$B$} \\
    \midrule
    IS & $\displaystyle Q\frac{\sinh\kap(T-t)}{\sinh\kap T}$
       & $y(0)=Q,\; y(T)=0$
       & $-Q\coth\kap T$ \\[8pt]
    TC & $\displaystyle Q\!\left(1 - \frac{\sinh\kap t}{\sinh\kap T}\right)$
       & $y(0)=0,\; y(T)=-Q$
       & $\displaystyle -\frac{Q}{\sinh\kap T}$ \\[8pt]
    TWAP & $\displaystyle Q\!\left(1 - \frac{t}{T}\right)$
         & $y(0)=0,\; y(T)=0$
         & $0$ \\
    \bottomrule
  \end{tabular}
\end{center}

\medskip
The stochastic risk of each strategy, identified from the payoff via
integration by parts and It\^{o}'s isometry, is:
\[
  \Var(\Pi)
  = \sigma^2\!\int_0^T\!\bigl(q_t - f(t)\bigr)^2\dd t,
  \quad\text{where }
  f(t) = \begin{cases} 0 & \text{(IS)} \\ Q & \text{(TC)} \\ Q(T-t)/T & \text{(TWAP)} \end{cases}
\]
\end{resultbox}

\begin{figure}[htbp]
  \centering
  \includegraphics[width=0.78\linewidth]{\detokenize{Graphs/Optimal liquidation strategies.png}}
  \caption{Closed-form optimal inventory trajectories for IS, TC, and TWAP under several values of risk aversion $\lambda$ and temporary impact $\eta$. IS becomes more front-loaded as risk aversion rises or trading becomes cheaper, TC becomes more back-loaded, and TWAP remains exactly linear.}
  \label{fig:optimal-strategies-lambda-eta}
\end{figure}

% -------------------------------------------------------------------
\subsection{Plot the efficient frontiers}
% -------------------------------------------------------------------

\subsubsection{Implementation Shortfall}

We compute the IS cost and risk so we can place the schedule on the efficient frontier.

This gives the IS cost:

\begin{resultbox}[IS - closed-form cost]
\[
  \mathrm{Cost}_{\mathrm{IS}}
  = \frac{\eta Q^2 \kap}{2}
    \left(\coth(\kap T) + \frac{\kap T}{\sinh^2(\kap T)}\right).
\]
\end{resultbox}

This gives the IS risk:

\begin{resultbox}[IS - closed-form risk]
\[
  \mathrm{Risk}_{\mathrm{IS}}
  = \frac{\sigma^2 Q^2}{2\kap}
    \left(\coth(\kap T) - \frac{\kap T}{\sinh^2(\kap T)}\right).
\]
\end{resultbox}

% -------------------------------------------------------------------
\subsubsection{Target Close}
% -------------------------------------------------------------------

We compute the TC cost and risk so we can compare its frontier with IS.

\paragraph{Cost.}
We start from the optimal TC trading speed:
\[
  v_{\mathrm{TC}}(t) = Q\,\frac{\kap\cosh(\kap t)}{\sinh(\kap T)}.
\]
This gives
\begin{align*}
  \mathrm{Cost}_{\mathrm{TC}}
  &= \eta\int_0^T v_{\mathrm{TC}}(t)^2\,\dd t
   = \eta Q^2\,\frac{\kap^2}{\sinh^2(\kap T)}\int_0^T \cosh^2(\kap t)\,\dd t.
\end{align*}
The integral $\int_0^T \cosh^2(\kap t)\,\dd t$ is the same one as in the IS case:
\[
  \int_0^T \cosh^2(\kap t)\,\dd t
  = \frac{T}{2} + \frac{\sinh(\kap T)\cosh(\kap T)}{2\kap}.
\]
This gives the TC cost:

\begin{resultbox}[TC - closed-form cost]
\[
  \mathrm{Cost}_{\mathrm{TC}}
  = \frac{\eta Q^2 \kap}{2}
    \left(\coth(\kap T) + \frac{\kap T}{\sinh^2(\kap T)}\right)
  = \mathrm{Cost}_{\mathrm{IS}}.
\]
\end{resultbox}

\paragraph{Risk.}
For TC, the risk term uses $(q_t - Q)^2$.
We use
\[
  q_{\mathrm{TC}}(t) - Q = -Q\,\frac{\sinh(\kap t)}{\sinh(\kap T)}.
\]
This gives
\begin{align*}
  \mathrm{Risk}_{\mathrm{TC}}
  &= \sigma^2\int_0^T \bigl(q_{\mathrm{TC}}(t) - Q\bigr)^2\,\dd t
   = \sigma^2 Q^2\,\frac{1}{\sinh^2(\kap T)}\int_0^T \sinh^2(\kap t)\,\dd t.
\end{align*}
The integral $\int_0^T \sinh^2(\kap t)\,\dd t$ is the same one as in the IS case:
\[
  \int_0^T \sinh^2(\kap t)\,\dd t
  = \frac{\sinh(\kap T)\cosh(\kap T)}{2\kap} - \frac{T}{2}.
\]
This gives the TC risk:

\begin{resultbox}[TC - closed-form risk]
\[
  \mathrm{Risk}_{\mathrm{TC}}
  = \frac{\sigma^2 Q^2}{2\kap}
    \left(\coth(\kap T) - \frac{\kap T}{\sinh^2(\kap T)}\right)
  = \mathrm{Risk}_{\mathrm{IS}}.
\]
\end{resultbox}

% -------------------------------------------------------------------
\paragraph{Why the frontiers coincide.}
% -------------------------------------------------------------------

The equality of cost and risk for IS and TC is not a coincidence.
The two schedules are related by time reversal:
\[
  q_{\mathrm{TC}}(t) = Q - q_{\mathrm{IS}}(T - t).
\]
Differentiating gives $v_{\mathrm{TC}}(t) = v_{\mathrm{IS}}(T-t)$.
So the cost terms are the same after the change of variable $t \mapsto T - t$.
For the risk,
\[
  q_{\mathrm{TC}}(t) - Q = -q_{\mathrm{IS}}(T-t),
\]
so the risk terms are also the same after the same time reversal.
This gives the common efficient frontier:

\begin{resultbox}[IS and TC - identical efficient frontiers]
The IS and TC efficient frontiers coincide exactly:
\[
  \mathrm{Cost}^\star(\kap) = \frac{\eta Q^2 \kap}{2}
    \left(\coth(\kap T) + \frac{\kap T}{\sinh^2(\kap T)}\right)
    \]
  
\[
  \mathrm{Risk}^\star(\kap) = \frac{\sigma^2 Q^2}{2\kap}
    \left(\coth(\kap T) - \frac{\kap T}{\sinh^2(\kap T)}\right),
\]
for all $\kap = \sqrt{\lambda\sigma^2/\eta} > 0$.
\end{resultbox}

\begin{figure}[htbp]
  \centering
  \includegraphics[width=\linewidth]{\detokenize{Graphs/Efficient frontiers — IS vs TC vs TWAP.png}}
  \caption{Efficient frontiers obtained by sweeping $\lambda$: moving along the frontier traces the cost-variance trade-off, while TWAP appears as a single benchmark-aligned point because its optimal schedule is independent of $\lambda$.}
  \label{fig:efficient-frontiers}
\end{figure}

% -------------------------------------------------------------------
\subsection{Bonus: find the optimal strategy in the expected utility framework}
\label{sec:bonus-expected-utility}
% -------------------------------------------------------------------

We solve the bonus question in the expected-utility framework for the \emph{Target Close} (TC) and \emph{TWAP} benchmarks only.
We keep the same model and notation as above and do not repeat the cash-flow derivations.

\begin{keybox}[Key point: State variables]
  The two HJB problems do not have the same primitive state variables:
  \begin{itemize}[leftmargin=1.5em]
    \item for TC, the state is $(t,x,s,q)$ because the benchmark depends on the
          terminal price $S_T$;
    \item for TWAP, the state is $(t,x,s,a,q)$, where
          \[
            A_t \coloneqq \frac{1}{T}\int_0^t S_u \dd u,
            \qquad
            \dd A_t = \frac{S_t}{T}\dd t,
          \]
          because the benchmark depends on the whole price path.
  \end{itemize}
\end{keybox}

\subsubsection{Expected utility criterion}

We replace the mean--variance criterion with CARA utility so we can compare the two frameworks.

We replace the mean--variance objective by CARA utility:
\[
  U(x) = -e^{-\gamma x},
  \qquad \gamma > 0.
\]
For a benchmark-dependent payoff $\Pi_B$, the optimisation problem is
\[
  \max_q \E[-e^{-\gamma \Pi_B}],
  \qquad
  B \in \{\mathrm{TC}, \mathrm{TWAP}\}.
\]
For IS, we do not need a separate HJB derivation.
Under a deterministic strategy, $\Pi_{\mathrm{IS}} = X_T - QS_0$ is Gaussian and has the same mean--variance structure as above.
The CARA--Gaussian identity therefore reduces the expected-utility problem to the same deterministic objective after the identification
$\lambda = \gamma/2$, so the optimal schedule is again
$q_{\mathrm{IS}}^*(t)=Q\,\sinh\!\bigl(\kap(T-t)\bigr)/\sinh(\kap T)$ with
$\kap = \sqrt{\gamma\sigma^2/(2\eta)}$.

% -------------------------------------------------------------------
\subsubsection{Target Close (TC)}
% -------------------------------------------------------------------

We solve the TC expected-utility problem so we can recover the optimal schedule from the HJB.

\paragraph{Value function and HJB.}
For the TC benchmark, the payoff is
\[
  \Pi_{\mathrm{TC}} = X_T - QS_T.
\]
The value function is
\begin{equation}
  V^{\mathrm{TC}}(t,x,s,q)
  =
  \sup_{(v_u)_{u \in [t,T]}}
  \E_{t,x,s,q}\!\left[-e^{-\gamma(X_T - QS_T)}\right].
  \label{eq:V-TC-bonus}
\end{equation}
Using the model dynamics,
\[
  \dd X_t = v_t(S_t-\eta v_t)\dd t,
  \qquad
  \dd S_t = \sigma \dd W_t,
  \qquad
  \dd q_t = -v_t \dd t,
\]
the HJB equation is
\begin{equation}
  0
  =
  V_t^{\mathrm{TC}}
  +
  \frac{1}{2}\sigma^2 V_{ss}^{\mathrm{TC}}
  +
  \sup_v
  \left\{
    v(s-\eta v)V_x^{\mathrm{TC}} - vV_q^{\mathrm{TC}}
  \right\}.
  \label{eq:HJB-TC-bonus}
\end{equation}

\paragraph{Exponential ansatz.}
Since $q_T = 0$, the terminal payoff can be rewritten as
\[
  X_T - QS_T = X_T + (q_T - Q)S_T.
\]
This motivates the ansatz
\begin{equation}
  V^{\mathrm{TC}}(t,x,s,q)
  =
  -\exp\!\left(
    -\gamma\bigl(x+(q-Q)s-\phi^{\mathrm{TC}}(t,q)\bigr)
  \right).
  \label{eq:ansatz-TC-bonus}
\end{equation}
The derivatives are
\[
  V_t^{\mathrm{TC}} = \gamma \phi_t^{\mathrm{TC}} V,
  \qquad
  V_x^{\mathrm{TC}} = -\gamma V,
\]
\[
  V_q^{\mathrm{TC}} = -\gamma(s-\phi_q^{\mathrm{TC}})V,
  \qquad
  V_s^{\mathrm{TC}} = -\gamma(q-Q)V,
\]
\[
  V_{ss}^{\mathrm{TC}} = \gamma^2(q-Q)^2V.
\]

\paragraph{Reduced HJB.}
Substituting these derivatives into \eqref{eq:HJB-TC-bonus} gives
\[
  \gamma \phi_t^{\mathrm{TC}} V
  +
  \frac{1}{2}\sigma^2\gamma^2(q-Q)^2V
  +
  \sup_v
  \left\{
    (vs-\eta v^2)(-\gamma V)
    -
    v\bigl(-\gamma(s-\phi_q^{\mathrm{TC}})V\bigr)
  \right\}
  = 0.
\]
Expanding the control term,
\[
  (vs-\eta v^2)(-\gamma V)
  -
  v\bigl(-\gamma(s-\phi_q^{\mathrm{TC}})V\bigr)
  =
  \gamma\eta v^2V - \gamma v\phi_q^{\mathrm{TC}}V.
\]
Since $-\gamma V > 0$, dividing through by $-\gamma V$ yields
\begin{equation}
  -\phi_t^{\mathrm{TC}}
  -
  \frac{\gamma\sigma^2}{2}(q-Q)^2
  +
  \sup_v
  \left\{
    v\phi_q^{\mathrm{TC}} - \eta v^2
  \right\}
  = 0.
  \label{eq:reduced-HJB-TC-bonus}
\end{equation}
The first-order condition is
\[
  \phi_q^{\mathrm{TC}} - 2\eta v = 0,
  \qquad
  v_{\mathrm{TC}}^*(t,q) = \frac{\phi_q^{\mathrm{TC}}(t,q)}{2\eta}.
\]
So the reduced HJB can be written as
\begin{equation}
  \phi_t^{\mathrm{TC}}
  +
  \inf_v
  \left\{
    \eta v^2 - v\phi_q^{\mathrm{TC}} + \frac{\gamma\sigma^2}{2}(q-Q)^2
  \right\}
  = 0.
  \label{eq:TC-running-cost-bonus}
\end{equation}
So the cost at each instant is
\[
  \eta v_t^2 + \frac{\gamma\sigma^2}{2}(q_t-Q)^2.
\]
\begin{equation}
  J_{\mathrm{TC}}^{\mathrm{EU}}[q]
  =
  \eta\!\int_0^T \dot q_t^2 \dd t
  +
  \frac{\gamma\sigma^2}{2}\!\int_0^T (q_t-Q)^2 \dd t.
  \label{eq:J-TC-bonus}
\end{equation}
Since the strategy is deterministic, $\Pi_{\mathrm{TC}}$ is Gaussian.
The CARA--Gaussian identity then maps the expected-utility problem to the mean--variance objective above with $\lambda = \gamma/2$.

\paragraph{Optimal TC strategy.}
The Euler--Lagrange equation associated with \eqref{eq:J-TC-bonus} is
\[
  \frac{\dd}{\dd t}(2\eta \dot q_t) - \gamma\sigma^2(q_t-Q) = 0,
\]
that is,
\begin{equation}
  \ddot q_t = \kap^2(q_t-Q),
  \qquad
  \kap^2 = \frac{\gamma\sigma^2}{2\eta}.
  \label{eq:EL-TC-bonus}
\end{equation}
We write $y_t \coloneqq q_t-Q$, where $y_t$ is the gap from a full position.
Then
\[
  \ddot y_t = \kap^2 y_t,
  \qquad
  y_0 = 0,
  \qquad
  y_T = -Q.
\]
The general solution is
\[
  y_t = A\sinh(\kap t) + B\cosh(\kap t).
\]
Using $y_0 = 0$ gives $B=0$.
Using $y_T = -Q$ gives
\[
  A = -\frac{Q}{\sinh(\kap T)}.
\]
So
\[
  y_t = -Q\frac{\sinh(\kap t)}{\sinh(\kap T)},
\]
and the optimal TC inventory is
\begin{equation}
  \boxed{
    q_{\mathrm{TC}}^*(t)
    =
    Q\left(1-\frac{\sinh(\kap t)}{\sinh(\kap T)}\right).
  }
  \label{eq:qstar-TC-bonus}
\end{equation}

% -------------------------------------------------------------------
\subsubsection{TWAP}
% -------------------------------------------------------------------

We solve the TWAP expected-utility problem so we can check whether the linear schedule remains optimal.

\paragraph{Value function and HJB.}
For TWAP, the benchmark is the running time average
\[
  A_T = \frac{1}{T}\int_0^T S_u \dd u,
\]
so the payoff is
\[
  \Pi_{\mathrm{TWAP}} = X_T - QA_T.
\]
The value function is
\begin{equation}
  V^{\mathrm{TWAP}}(t,x,s,a,q)
  =
  \sup_{(v_u)_{u \in [t,T]}}
  \E_{t,x,s,a,q}\!\left[-e^{-\gamma(X_T-QA_T)}\right].
  \label{eq:V-TWAP-bonus}
\end{equation}
The HJB equation now includes the extra state variable $a=A_t$:
\begin{equation}
  0
  =
  V_t^{\mathrm{TWAP}}
  +
  \frac{1}{2}\sigma^2 V_{ss}^{\mathrm{TWAP}}
  +
  \frac{s}{T}V_a^{\mathrm{TWAP}}
  +
  \sup_v
  \left\{
    v(s-\eta v)V_x^{\mathrm{TWAP}} - vV_q^{\mathrm{TWAP}}
  \right\}.
  \label{eq:HJB-TWAP-bonus}
\end{equation}

\paragraph{Exponential ansatz.}
To match the terminal payoff $X_T-QA_T$, we use
\begin{equation}
  V^{\mathrm{TWAP}}(t,x,s,a,q)
  =
  -\exp\!\left(
    -\gamma\left(
      x+\left(q-Q\frac{T-t}{T}\right)s-Qa-\phi^{\mathrm{TWAP}}(t,q)
    \right)
  \right).
  \label{eq:ansatz-TWAP-bonus}
\end{equation}
At time $T$, the coefficient of $s$ vanishes, so this becomes
$x-Qa$.

From \eqref{eq:ansatz-TWAP-bonus}, the derivatives are
\[
  V_t^{\mathrm{TWAP}}
  =
  \gamma\left(\phi_t^{\mathrm{TWAP}}-\frac{Q}{T}s\right)V,
\]
\[
  V_x^{\mathrm{TWAP}} = -\gamma V,
  \qquad
  V_q^{\mathrm{TWAP}} = -\gamma(s-\phi_q^{\mathrm{TWAP}})V,
\]
\[
  V_s^{\mathrm{TWAP}}
  =
  -\gamma\left(q-Q\frac{T-t}{T}\right)V,
  \qquad
  V_{ss}^{\mathrm{TWAP}}
  =
  \gamma^2\left(q-Q\frac{T-t}{T}\right)^2V,
\]
\[
  V_a^{\mathrm{TWAP}} = \gamma QV.
\]

\paragraph{Reduced HJB.}
Substituting these derivatives into \eqref{eq:HJB-TWAP-bonus}, the terms
involving $Qs/T$ cancel. Dividing through by $-\gamma V > 0$ yields
\begin{equation}
  -\phi_t^{\mathrm{TWAP}}
  -
  \frac{\gamma\sigma^2}{2}
  \left(q-Q\frac{T-t}{T}\right)^2
  +
  \sup_v
  \left\{
    v\phi_q^{\mathrm{TWAP}} - \eta v^2
  \right\}
  = 0.
  \label{eq:reduced-HJB-TWAP-bonus}
\end{equation}
Equivalently,
\begin{equation}
  \phi_t^{\mathrm{TWAP}}
  +
  \inf_v
  \left\{
    \eta v^2
    -
    v\phi_q^{\mathrm{TWAP}}
    +
    \frac{\gamma\sigma^2}{2}
    \left(q-Q\frac{T-t}{T}\right)^2
  \right\}
  = 0.
  \label{eq:TWAP-running-cost-bonus}
\end{equation}
So the cost at each instant is
\[
  \eta v_t^2
  +
  \frac{\gamma\sigma^2}{2}
  \left(q_t-Q\frac{T-t}{T}\right)^2.
\]
\begin{equation}
  J_{\mathrm{TWAP}}^{\mathrm{EU}}[q]
  =
  \eta\!\int_0^T \dot q_t^2 \dd t
  +
  \frac{\gamma\sigma^2}{2}\!\int_0^T
  \left(q_t-Q\frac{T-t}{T}\right)^2 \dd t.
  \label{eq:J-TWAP-bonus}
\end{equation}
Since the strategy is deterministic, $\Pi_{\mathrm{TWAP}}$ is Gaussian.
The CARA--Gaussian identity then maps the expected-utility problem to the mean--variance objective above with $\lambda = \gamma/2$.

\paragraph{Optimal TWAP strategy.}
The Euler--Lagrange equation associated with \eqref{eq:J-TWAP-bonus} is
\begin{equation}
  \ddot q_t
  =
  \kap^2\left(q_t-Q\frac{T-t}{T}\right).
  \label{eq:EL-TWAP-bonus}
\end{equation}
We write
\[
  y_t \coloneqq q_t-Q\frac{T-t}{T}.
\]
where $y_t$ is the deviation from the TWAP benchmark path.
Since $Q(T-t)/T$ is linear in $t$, its second derivative is zero, so
\[
  \ddot y_t = \kap^2 y_t.
\]
The boundary conditions become
\[
  y_0 = q_0-Q = 0,
  \qquad
  y_T = q_T-0 = 0.
\]
The only solution is therefore
\[
  y_t \equiv 0.
\]
So the optimal TWAP inventory is
\begin{equation}
  \boxed{
    q_{\mathrm{TWAP}}^*(t)
    =
    Q\frac{T-t}{T}.
  }
  \label{eq:qstar-TWAP-bonus}
\end{equation}

\paragraph{Comparison with the mean--variance framework.}

This gives the final expected-utility result:

\begin{resultbox}[Expected utility - final answer]
  The expected-utility and mean--variance formulations coincide under
  \[
    \boxed{\lambda = \frac{\gamma}{2}}.
  \]
  The optimal strategies are
  \[
    q_{\mathrm{IS}}^*(t)
    =
    Q\,\frac{\sinh\!\bigl(\kap(T-t)\bigr)}{\sinh(\kap T)},
    \qquad
    q_{\mathrm{TC}}^*(t)
    =
    Q\left(1-\frac{\sinh(\kap t)}{\sinh(\kap T)}\right),
    \qquad
    q_{\mathrm{TWAP}}^*(t)
    =
    Q\frac{T-t}{T},
  \]
  with
  \[
    \kap = \sqrt{\frac{\gamma\sigma^2}{2\eta}}.
  \]
\end{resultbox}

% -------------------------------------------------------------------
\subsection{Bellman Equation on a Discretized Grid}
% -------------------------------------------------------------------

We also solve the IS and TC problems by dynamic programming on a discrete state space so we can compare them with the continuous-time formulas.
We split the horizon $[0,T]$ into $N$ intervals of length $\Delta t = T/N$, and we restrict the inventory to the grid
\[
  q_i = i\,\Delta q,
  \qquad
  \Delta q = \frac{Q}{n_q},
  \qquad
  i = 0,1,\dots,n_q.
\]
At each time step $k$ and current inventory $q_i$, the control is the number of shares sold over $[t_k,t_{k+1}]$.
Equivalently, we choose the next inventory level $q_j \le q_i$, so the trade size is $n_k = q_i-q_j$ and the selling rate is approximated by $v_k = (q_i-q_j)/\Delta t$.

We define the value function $V_k(q_i)$ as the minimal remaining cost from time $t_k$ onward when the current inventory is $q_i$.
In both benchmarks, we enforce the terminal liquidation constraint exactly through
\[
  V_N(0)=0,
  \qquad
  V_N(q_i)=+\infty \quad \text{for } q_i>0,
\]
which forces the strategy to end with zero inventory.
We then run a standard backward induction.
At each node $(t_k,q_i)$, we test all feasible next inventories $q_j \le q_i$, store the minimiser, and recover the optimal inventory path with a forward pass from $q_0 = Q$.

\subsubsection{Implementation Shortfall}

We write the discrete IS recursion so we can compare it with the closed-form solution.

For the implementation shortfall benchmark, the discrete Bellman equation is
\begin{equation}
  V_k^{\mathrm{IS}}(q_i)
  =
  \min_{0 \le j \le i}
  \left\{
    \eta\left(\frac{q_i-q_j}{\Delta t}\right)^2 \Delta t
    +
    \lambda\sigma^2 q_i^2 \Delta t
    +
    V_{k+1}^{\mathrm{IS}}(q_j)
  \right\}.
  \label{eq:bellman-is-discrete}
\end{equation}
The first term is the temporary impact cost over the current interval.
The second term is the discrete version of the IS risk penalty $\lambda\sigma^2\int_0^T q_t^2\,\dd t$.
We compare this numerical solution with the closed-form inventory path from Question~1.1,
\[
  q_{\mathrm{IS}}^*(t)
  =
  Q\,\frac{\sinh\!\bigl(\kap(T-t)\bigr)}{\sinh(\kap T)},
  \qquad
  \kap = \sqrt{\frac{\lambda\sigma^2}{\eta}}.
\]

In the notebook, we use $Q=1$, $T=1$, $\sigma=0.5$, $\eta=0.01$, $\lambda=1$, and three grid resolutions: $(N,n_q)=(5,10)$, $(10,20)$, and $(60,120)$.
Figure~\ref{fig:bellman-is} shows that the coarsest grids produce a visibly quantized inventory path, while the refined grid is almost indistinguishable from the analytical schedule.
This is the expected convergence pattern.
As $\Delta t \to 0$ and $\Delta q \to 0$, the discrete control set becomes dense in the continuous one and the Bellman solution recovers the front-loaded IS rule.

\begin{figure}[htbp]
  \centering
  \includegraphics[width=\linewidth]{\detokenize{Graphs/IS — Bellman vs analytical'.png}}
  \caption{Implementation Shortfall inventory paths. The Bellman solution is plotted for increasingly fine time and inventory grids and compared with the closed-form Almgren--Chriss schedule. Grid refinement makes the discrete trajectory converge to the analytical IS solution.}
  \label{fig:bellman-is}
\end{figure}

\subsubsection{Target Close}

We write the discrete TC recursion so we can compare it with the closed-form solution.

For the target close benchmark, the Bellman recursion has the same structure, but the variance penalty is centred around the initial inventory $Q$ rather than around $0$:
\begin{equation}
  V_k^{\mathrm{TC}}(q_i)
  =
  \min_{0 \le j \le i}
  \left\{
    \eta\left(\frac{q_i-q_j}{\Delta t}\right)^2 \Delta t
    +
    \lambda\sigma^2 (q_i-Q)^2 \Delta t
    +
    V_{k+1}^{\mathrm{TC}}(q_j)
  \right\}.
  \label{eq:bellman-tc-discrete}
\end{equation}
This is the discrete version of the TC mean--variance criterion derived earlier.
We compare it with the continuous-time optimum
\[
  q_{\mathrm{TC}}^*(t)
  =
  Q\left(1-\frac{\sinh(\kap t)}{\sinh(\kap T)}\right),
  \qquad
  \kap = \sqrt{\frac{\lambda\sigma^2}{\eta}}.
\]

We use the same sequence of grids in the numerical experiment.
Figure~\ref{fig:bellman-tc} again shows that coarse discretizations produce a staircase approximation, while the fine grid closely tracks the theoretical TC inventory path.
The qualitative difference with IS remains clear: the TC strategy is back-loaded, with little trading at the beginning and rapid liquidation near the close.
The near overlap between the fine-grid Bellman curve and the analytical one confirms that the discretized dynamic-programming scheme is consistent with the continuous-time solution.

\begin{figure}[htbp]
  \centering
  \includegraphics[width=\linewidth]{\detokenize{Graphs/TC — Bellman vs analytical.png}}
  \caption{Target Close inventory paths. The discrete Bellman solutions obtained on finer grids progressively approach the closed-form TC schedule, confirming convergence of the numerical approximation.}
  \label{fig:bellman-tc}
\end{figure}

\paragraph{Convergence to the theoretical solution.}
For both IS and TC, the convergence is clear from the inventory paths.
On coarse grids, there are only a few possible trades, so the optimal path is necessarily quantized.
As we refine the discretization, the Bellman recursion has access to a richer set of feasible controls and the discrete optimum converges to the hyperbolic-sine solutions from Question~1.1.
In particular, on the finest grid used in the notebook, the Bellman and analytical curves are nearly indistinguishable in both benchmarks.

\end{document}
